\documentclass[../main]{subfiles}
\begin{document}
\ifSubfilesClassLoaded{\setcounter{page}{4}}{}
\section{I fondamenti del comunismo}
\label{sec:03_fondamenti}

I fondamenti del comunismo risiedono nella produzione di merci. Nell'ambito del sistema capitalista, l'intera ricchezza sociale si concretizza in un insieme di beni. Poiché, nel capitalismo, i prodotti del lavoro assumono ineluttabilmente una forma commerciale, la merce rappresenta la "cellula" di tale sistema.

I principi del comunismo, inteso come sistema sociale, e in accordo con la concezione materialistica della storia, possono essere identificati in un corrispondente modo di produzione.

La ricchezza intrinseca della società comunista risiede nell’universalità dei contenuti e delle forme dell’attività umana. Il marxismo, in linea generale, concepisce l’essenza dell’uomo come un’attività di trasformazione della natura, compresa la sua stessa natura sociale. Il lavoro, la manipolazione della materia prima, la creazione di oggetti che non preesistono in natura: ecco ciò che distingue radicalmente l’esistenza umana da quella degli animali e che determina tutte le altre differenze. Ma fino ad oggi, l’attività lavorativa su scala sociale è stata priva di libertà. Tutti i sistemi di produzione precedenti al comunismo si sono contraddistinti per una specifica forma di lavoro non libero. Il comunismo, al contrario, postula il superamento totale di questa mancanza di libertà. Di conseguenza, la sua base può essere solo il lavoro libero, «privo di vincoli e di retribuzione», come «prima necessità dell’esistenza» umana.

La forma più elementare di attività (o, più precisamente, di impegno, poiché il termine "lavoro" conserva in qualche misura una connotazione negativa) è quella "cellula" del comunismo, nella quale risiede il germe dello sviluppo futuro della società comunista in tutte le sue necessarie manifestazioni.

L'attività libera e la sua intrinseca dimensione collettiva rappresentano una realtà tangibile, al pari di quanto lo fossero, prima dell'avvento del capitalismo, la produzione di beni e il concetto stesso di merce. La transizione verso un nuovo sistema sociale si compirà nel momento in cui l'attività libera dell'uomo si affermerà come principio universale, analogamente a quanto accaduto in passato con la merce. Dalla molteplicità di beni presenti sul mercato, ognuno caratterizzato da peculiarità specifiche, non è possibile desumere immediatamente la natura della merce e il processo produttivo ad essa legato. Allo stesso modo, dalla diversità di forme e modalità di organizzazione dell'attività comunista, non si può cogliere appieno la sua essenza e la sua forma intrinsecamente collettiva. Inoltre, le parole possono talvolta velare la sostanza delle cose. La forma immediatamente collettiva dell'attività libera può manifestarsi anche qualora un singolo individuo la persegua in solitudine.È opportuno sottolineare le definizioni che risultano cruciali per la sua comprensione.

Nella sua forma più compiuta e collettiva, l’attività libera si manifesta in conformità ai principi di libertà, verità, bene e bellezza, incarnando così anche l’anelito più profondo dell’animo umano. Si tratta, in definitiva, di una forma di espressione sensoriale intrinsecamente preziosa. L’attività libera implica, per sua natura, l’armoniosa integrazione di sensibilità, moralità e intelletto, elevandoli congiuntamente al livello della piena libertà. In questo processo, l’individuo si afferma nella sua totalità, trascendendo la natura specifica dell’azione compiuta e l’eventuale preminenza di uno specifico valore – verità, bene o bellezza.

Questa peculiare modalità, intrinsecamente collettiva, rappresenta la struttura stessa dell’azione, osservata dall’esterno e considerata in relazione al suo contenuto. Nel linguaggio corrente, tale attività è spesso definita creatività sociale, o più semplicemente, creatività. La forma riveste un’importanza cruciale, poiché lo sviluppo si manifesta e si esprime attraverso il dinamico mutare delle forme. Il contenuto, inevitabilmente, dà vita a una forma corrispondente, e si concretizza proprio attraverso di essa. Di conseguenza, le contraddizioni insite nel contenuto vengono inizialmente percepite dalla coscienza come contraddizioni formali. È attraverso il suo dispiegarsi, e nel dispiegarsi stesso (la trasformazione delle forme), che si compie effettivamente il movimento del contenuto. Dal punto di vista strutturale, nell’ambito di questa attività, elementi quali il tipo di relazione tra gli individui, l’obiettivo perseguito, l’esistenza sociale dell’oggetto e dei mezzi impiegati, assumono un ruolo di primo piano.

L'attività collettiva, intesa come espressione di una comunità orientata a obiettivi comuni, si manifesta primariamente nella sua attenzione verso gli altri. In queste forme di collaborazione, l'individuo agisce sempre perseguendo il bene comune, elevandolo a fine ultimo e valore supremo. Ogni azione è diretta allo sviluppo e al benessere degli altri. Di conseguenza, il prodotto di tale impegno non assume mai una forma autonoma o puramente materiale, ma piuttosto un valore intrinseco. Esso non si pone come un'entità estranea e ostile nei confronti dei membri della comunità, bensì come elemento integrante e funzionale al loro progresso. In sintesi, questa attività è volta a favorire lo sviluppo sociale e a migliorare la qualità della vita di tutti i suoi partecipanti, realizzandosi così anche per i singoli, in quanto parte integrante di tale società. Ciò che oggi viene definito come "valori materiali" nell'ambito dell'attività comunitaria può essere considerato unicamente come un mezzo a tale fine.

Per quanto concerne la natura delle relazioni interpersonali, queste possono essere genericamente definite come rapporti di cameratismo, o, in termini più ampi, come forme di associazione libera.

Questo tipo di interazione rivela altresì le modalità attraverso cui si svolge l'attività, poiché in un'associazione fondata sulla libera volontà, nessuno dei partecipanti può essere ridotto a mero strumento per l'attuazione delle aspirazioni e degli interessi altrui. Sebbene un individuo, grazie alle proprie doti e competenze, possa assumere un ruolo di guida all'interno del gruppo, gli altri membri non si limitano a eseguire passivamente le sue direttive, bensì contribuiscono attivamente al perseguimento di un obiettivo condiviso, ugualmente rilevante e vantaggioso per tutti.

Questa forma di relazione evidenzia inoltre la natura intrinsecamente sociale dei mezzi e degli scopi dell'azione comunista. Pertanto, solo elementi o processi distinti dagli esseri umani possono essere impiegati come strumenti. La dimensione sociale dell'esistenza dei mezzi e degli oggetti del lavoro implica il loro utilizzo in un contesto di interazione diretta. È in questa prospettiva che Marx teorizzò il passaggio da una gestione incentrata sulle persone a una gestione orientata alle cose.

Come si evince da quanto precede, tali forme di attività umana sono esenti da qualsiasi tipo di costrizione esterna, sia essa diretta o indiretta (come nel caso della retribuzione per il lavoro), poiché il processo e il risultato dell'attività comunista non sono affatto indifferenti per chi la compie. In ogni caso, l'individuo è direttamente coinvolto nel proprio lavoro, a differenza di quanto accade nella società moderna, dove l'attività stessa si riduce a mero strumento per procurarsi denaro.

È innegabile che la vera anima di una comunità risieda nella sua interdipendenza: l'attività di ciascun individuo è intrinsecamente legata al contributo degli altri, poiché gli strumenti, le risorse e le metodologie impiegate sono il frutto del lavoro altrui. A sua volta, l'impegno di ogni membro si integra nel tessuto del lavoro collettivo, offrendo al contempo il materiale per le attività future. Tuttavia, questo principio si applica a qualsiasi forma di lavoro e organizzazione sociale, e non costituisce, di per sé, un elemento distintivo del comunismo.

Per attività collettiva diretta intendiamo un’unione di individui impegnati in un processo collaborativo, in cui lo sviluppo di ciascuno è finalizzato al progresso dell’intero gruppo, e i mezzi per conseguire tale scopo sono costituiti da elementi e processi naturali, in relazione ai quali si manifesta una volontà libera e su un piano di parità.

Questa forma di attività collettiva non rappresenta semplicemente un mero involucro sociale, un’organizzazione esterna che può coesistere o meno senza incidere sul contenuto dell’azione stessa; al contrario, essa ne costituisce l’essenza. Al di fuori di questa struttura, non è possibile realizzare un’attività autenticamente libera, fondata sui principi di verità, bene e bellezza.\footnote{Contrariamente a quanto si potrebbe pensare, ciò non si limita alla semplice associazione tra bellezza e gusto, tra bene e opinioni divergenti sul giusto e sullo sbagliato; al contrario, si tratta di un approccio che attribuisce un significato oggettivo alla verità.Nel corso dei millenni, le categorie di «verità», «bene» e «bellezza» sono state definite come attributi ideali e oggettivi dell’agire umano nel mondo, manifestandosi – o meglio, evolvendosi – all’interno di tale agire e, per estensione, configurandosi anche come proprietà oggettive del mondo stesso (della materia). In sintesi, quando si fa riferimento a un’azione valutata in base a tali criteri, ci si riferisce ai «momenti di progresso» della sensibilità, della conoscenza e della moralità umane, all’evoluzione delle relazioni interpersonali e del rapporto tra l’uomo e l’ambiente, considerati nella loro dimensione ideale.} L’evoluzione e le trasformazioni di questa forma implicano la preservazione dei principi sopracitati, in modo analogo a quanto accade nelle metamorfosi della forma mercantile dei prodotti del lavoro, nelle quali il valore viene costantemente mantenuto (l’elemento che distingue un prodotto come bene di scambio, e non come semplice risultato di un’attività lavorativa).

L’essenza di tale attività risiede nella sua intrinseca natura rivoluzionaria. Ed è proprio in questa qualità che si annida la sua contraddizione interna. Nell’ambito dell’attività comunista, liberata da vincoli, si definiscono costantemente sia i suoi confini intrinseci sia le possibilità di superarli. Attraverso le conquiste storiche – e dunque attraverso i limiti imposti dalla storia – l’attività stessa subisce una trasformazione, arricchendosi di nuovi contenuti (grazie all’estensione del campo di dominio della materia) e assumendo nuove forme. Questa evoluzione non avviene in maniera indipendente dai soggetti coinvolti nel processo lavorativo, come accade invece nel processo di produzione materiale. Al contrario, tale trasformazione si compie in modo consapevole e mirato. Pertanto, il comunismo, inteso come dominio dell’uomo sulla propria attività, svela la contraddizione intrinseca all’attività umana in generale, e di conseguenza la contraddizione insita nel mondo materiale in cui l’uomo esercita la propria azione.Per questo motivo, anche nell’ambito della scienza comunista, l’osservazione di Marx espressa nella prima tesi su Feuerbach riveste un’importanza metodologica capitale: «Il limite principale di tutto il materialismo precedente, incluso quello di Feuerbach, risiede nel fatto che l’oggetto, la realtà, la sensibilità vengono concepiti unicamente come oggetti, o come forme di contemplazione, e non come attività sensoriale umana, come prassi, intesa in senso oggettivo. Ne consegue che, in opposizione al materialismo, si è sviluppata una concezione attiva, ma essa rimane confinata in una dimensione astratta, poiché l’idealismo, per sua natura, ignora la vera e propria attività sensoriale. Feuerbach intende confrontarsi con oggetti sensibili, concretamente distinti dalle idee di ciò che è buono e ciò che è cattivo, le quali possono variare, e solo in questo modo attribuisce un significato oggettivo alla verità.Le categorie di “verità”, “bene” e “bellezza” sono state progressivamente definite nel corso dei secoli come i tratti ideali e oggettivi che caratterizzano l’attività umana nel mondo, manifestandosi – o, meglio, evolvendosi – all’interno di tale attività e costituendo, perciò, anche attributi oggettivi del mondo stesso (della materia). In sintesi, quando si analizza l’attività umana alla luce di questi concetti – verità, bene e bellezza – ci si riferisce ai “punti di convergenza” della sensibilità, della conoscenza e della moralità umana, allo sviluppo delle relazioni tra gli individui e tra l’uomo e la natura, il tutto considerato da una prospettiva ideale. È importante notare che, pur avendo come riferimento degli oggetti, l’attività umana non viene qui considerata in funzione di uno scopo specifico. In questa prospettiva, si evidenzia il ruolo cruciale dell’idealismo nella genesi della teoria dell’attività.Pur essendo un compito di tale portata e rilevanza, tanto che in seguito Lenin sottolineò l’importanza di formare una comunità di “adepti del materialismo dialettico”, le categorie della dialettica, intese in chiave materialistica come “forme di attività umana ripetute innumerevoli volte” (Lenin), la loro intrinseca contraddittorietà e la loro dinamicità rappresentano uno strumento indispensabile per l’elaborazione della teoria comunista, poiché esse riflettono i momenti essenziali dell’attività umana in generale, consentendo di individuare gli elementi senza i quali tale attività non potrebbe realizzarsi. Nella tesi di Marx, tuttavia, viene altresì posto l’accento sulla concretezza dell’attività umana.

È bene ribadire che l'attività libera non è un semplice «riposo» nel senso contemporaneo del termine, né un mero passatempo, bensì un'azione che si conforma ai dettami della libertà. In passato, la filosofia, per bocca di Baruch Spinoza, definì la libertà come una forma di necessità consapevole. La vera libertà, in definitiva, si manifesta attraverso un'azione coerente con tale necessità.

È nel tempo libero, e in funzione delle esigenze storiche, che gli individui coltivano la cultura in senso lato. Questa attività non è un’invenzione comunista, un’illusione nata dalla fantasia; al contrario, essa ha radici nel passato e persiste ancora oggi. Tuttavia, in precedenza come ora, la sua portata è limitata e riservata a una ristretta cerchia di individui. La stragrande maggioranza dei membri della società non ha accesso a questa forma di espressione libera. Questo deriva, in parte, dal fatto che il mero possesso di tempo libero non è sufficiente; è necessario che l’individuo possieda un certo grado di sviluppo personale per poterlo impiegare in attività significative. Le condizioni che favoriscono tale sviluppo sono state, e continuano ad essere, piuttosto ristrette all’interno della moderna società orientata alla produzione di merci.

Pertanto, per avanzare verso il comunismo, è imperativo promuovere la nascita del comunismo stesso, ovvero dell’attività libera; più precisamente, occorre aumentare il tempo libero a disposizione e, parallelamente, elevare il livello di sviluppo dei membri della società, affinché essi possano dedicarsi a tale attività.

Pertanto, la distribuzione a cui ci riferiamo è primariamente la distribuzione del tempo libero e delle opportunità per l'acquisizione della cultura umana. La distribuzione dei beni prodotti dal lavoro è necessaria solo nella misura in cui essa assicuri il tempo libero. In questa prospettiva, è fondamentale affrontare la questione dell'efficienza della produzione sociale durante il periodo di transizione verso il comunismo. Nel sistema capitalistico, il tempo libero rappresenta un costo di produzione, non viene preso in considerazione né valorizzato, e di conseguenza, spesso non si trasforma in un vero e proprio tempo libero, non venendo utilizzato per lo scopo che gli è proprio. L'obiettivo del tempo libero è, infatti, quello di essere impiegato per favorire il costante sviluppo delle relazioni sociali (o, più in generale, per lo sviluppo delle attività umane).

Per il comunismo, è di fondamentale importanza che l’idea – ormai diffusa e condivisa – di considerare il tempo libero e le attività che si svolgono in tale periodo come un bene comune, immune da utilizzi individuali e dispendiosi, sia ampiamente valorizzata. In generale, la forza dell’abitudine e la sua capacità di plasmare comportamenti collettivi rivestono un’enorme rilevanza, anche solo perché l’esercizio della libertà, ovvero l’agire in conformità ai suoi principi, si rivela sempre un compito arduo. Ciò implica confrontarsi costantemente con le contraddizioni insite nel processo di sviluppo sociale. Non si tratta soltanto di un impegno gravoso, ma anche di una sfida che richiede un superamento continuo dei propri limiti, una sorta di negazione del proprio passato. La libertà collettiva e sociale si configura, dunque, come una tragedia collettiva e sociale, e la sua perpetua risoluzione rappresenta la chiave per il progresso.La società comunista non è certo un sistema di indolenza e ozio generalizzati, bensì una società fondata sull’attività libera e universale, che costituisce per ciascun individuo una necessità primaria.

Il comunismo si concretizza attraverso una trasformazione intrinseca delle forme collettive dell’attività umana, mirando a rinnovare l’attività stessa, a vivificare l’esistenza di ogni persona. Le forme che si evolvono, nel tentativo di superare i propri limiti, danno origine a nuove possibilità. Pertanto, questo processo di cambiamento non può essere concepito come una mera ripetizione di metamorfosi successive (come si osserva nel dinamismo del capitale). La logica intrinseca al mutamento delle forme collettive dell’attività libera, la sua imprescindibile necessità, è determinata non solo dallo specifico modo di produzione storicamente raggiunto e dai rapporti di produzione ad esso inerenti, ma anche dalla libertà, che si manifesta come necessità intrinseca della materia, e che si rivela progressivamente con una comprensione sempre più profonda della stessa.

In un'accezione più ampia, il principio cardine del comunismo può essere espresso come legge della libertà, o legge dell'autodeterminazione dell'azione umana: l'attività libera dei membri della società, manifestandosi come un'evoluzione consapevole dei suoi contenuti e delle sue forme, determina il progresso dell'intera comunità. In altre parole, la libertà, una volta pienamente attuata, tende intrinsecamente a espandersi e rafforzarsi.

Naturalmente, la comprensione del funzionamento di tale principio richiede un continuo perfezionamento, avvalendosi di tutte le conoscenze disponibili, sia relative alla natura che alla società. Pertanto, la teoria del comunismo può, a determinate condizioni, assumere un ruolo unificante per tutte le discipline scientifiche, evidenziando il loro reale nesso, ovvero il legame stesso insito nella prassi umana.

Le attuali forme di attività umana, libere e collettive, definite come «germi del comunismo», non costituiscono una mera dimostrazione della sua fattibilità, bensì ne rappresentano la concreta esistenza. È fondamentale comprendere che il modello produttivo comunista è già presente, e la sfida consiste nel farlo diventare egemone e diffonderlo capillarmente nell’intera società.

Nelle loro fasi di sviluppo, le forme di attività umana intrinsecamente collettiva, da intendersi come i «semi del comunismo», possono evolvere fino a configurare un sistema comunista, ossia un modello di organizzazione sociale. In quest’ottica, uno dei pilastri più significativi dell’eredità teorica leniniana risiede nell’individuazione e nella promozione di tali «semi» del comunismo, ovunque essi si manifestino
\end{document}
